\scanpage{24}
\begin{center}- 80 -\end{center}

\uline{Définition 6.} \uline{Soient} $A$ \uline{un anneau d'intégrité}, $K$ \uline{son corps des fractions}
$P^*$ \uline{le groupe ordonné des idéaux principaux fractionnaires de $K$ relativement à} $A$; \uline{on dira que} $A$ \uline{est factoriel si} $P^*$ \uline{est isomorphe à un groupe} $\unsure{\mathbf Z|\mathbb Z}^{(I)}$.

En d'autres termes, $A$ est factoriel s'il existe une famille $(p_\iota)$ d'éléments de $A$ ayant la propriété suivante: tout $x\in K^*$ peut s'écrire d'une manière et d'une seule sous la forme
\[
x=u\prod_\iota p_\iota^{n_\iota}\qquad(1)
\]
où $u$ est un élément inversible de $A$, et où les $n_\iota$ sont des entiers, tous nuls sauf un nombre fini d'entre eux, et tous $\geq0$ chaque fois que $x\in A^*$. Le groupe $P^*$ étant alors réticulé, toutes les propriétés de divisibilité démontrées en algèbre (Alg.,Chap...) concernant le pgcd, le ppcm, les éléments étrangers, etc., s'appliquent aux anneaux factoriels; la famille $(p_\iota)$ étant comme ci-dessus, il est clair que tout $p_\iota$ est un ``élément extrémal'' de $A$ (Alg.,...), et que réciproquement tout élément extrémal est de la forme $up_\iota$, où $u$ est inversible dans $A$, c'est-à-dire est associé à un $p_\iota$ et à un seul.

Dans ces conditions, $A$ est normal. En effet, pour tout $x\in K^*$, soit $v_\iota(x)=n_\iota$ l'exposant de $p_\iota$ dans l'expression (1) de $x$; et soit $v_\iota(0)=+\infty$. D'après l'unicité de (1), on a $v_\iota(xy)=v_\iota(x)+v_\iota(y)$ quels que soient $x,y$. Par hypothèse, $x\in A$ entraîne $v_\iota(x)\geq0$ quel que soit $\iota$; la réciproque est évidente. Il s'ensuit que, pour $x\in A$, $x\in\unsure{p_\iota^{n_A}|p_\iota^n A}$ est équivalent à $v_\iota(x)\geq n$; on en conclut aussitôt que, pour $x\in A$, $v_\iota(x+y)\geq\min[v_\iota(x),v_\iota(y)]$; d'après la prop.2 du § 3, n°\unsure{2}, la même propriété reste valable alors quels que soient $x,y$ dans $K$; donc les $v_\iota$ sont des valuations discrètes de $A$ et forment un ensemble de définition de $A$, qui est donc normal. Le centre de $v_\iota$ est l'idéal principal $p_\iota A$; d'après le coroll.2 de la prop.4, n°2, c'est donc un
\margin{23}
